\subsection{For Count}

\begin{center}
    \texttt{Derivation Tree}
\end{center}

\begin{center}
\begin{forest}
for tree={
    font=\tiny,
    grow=south,
    parent anchor=south,
    child anchor=north,
    l sep=4mm,
    s sep=1mm,
    inner sep=1pt
}
[{$<$for\_count$>$}
    [{"for"}[{"for"},tier=word]]
    [{$<$space$>$}
        [{" "},tier=word]
    ]
    [{$<$expression$>$}
        [{$<$term$>$}
            [{$<$factor$>$}
                [{$<$literal$>$}
                    [{$<$number$>$}
                        [{$<$digit$>$}
                            [{"2"},tier=word]
                        ]
                    ]
                ]
            ]
        ]
    ]
    [{$<$space$>$}
        [{" "},tier=word]
    ]
    [{"times"}[{"times"},tier=word]]
    [{NEWLINE}
        [{"\textbackslash n"},tier=word]
    ]
    [{$<$block$>$}
        [{INDENT}
            [{"$\rightarrow$"},tier=word]
        ]
        [{$<$statement\_list$>$}
            [{$<$statement$>$}
                [{$<$assignment$>$}
                    [{$<$identifier$>$}
                        [{$<$letter$>$}
                            [{$<$lower$>$}
                                [{"x"},tier=word]
                            ]
                        ]
                    ]
                    [{$<$optional\_space$>$}[{" "},tier=word]]
                    [{"="}[{"="},tier=word]]
                    [{$<$optional\_space$>$}[{" "},tier=word]]
                    [{$<$expression$>$}
                        [{$<$term$>$}
                            [{$<$factor$>$}
                                [{$<$literal$>$}
                                    [{$<$number$>$}
                                        [{$<$digit$>$}
                                            [{"1"},tier=word]
                                        ]
                                    ]
                                ]
                            ]
                        ]
                    ]
                ]
                [{NEWLINE}[{"\textbackslash n"},tier=word]]
            ]
        ]
        [{DEDENT}
            [{"$\leftarrow$"},tier=word]
        ]
    ]
]
\end{forest}
\end{center}

\vspace{5mm}

\begin{center}
    \texttt{Simplified Abstract Syntax Tree}
\end{center}

\begin{center}
\begin{forest}
for tree={
    font=\small,
    grow=south,
    parent anchor=south,
    child anchor=north,
    l sep=6mm,
    s sep=4mm,
    inner sep=2pt
}
[{$<$for\_count$>$}
    [{"for"}[{"for"},tier=word]]
    [{$<$expression$>$}
        [{$<$number$>$}[{"2"},tier=word]]
    ]
    [{"times"}[{"times"},tier=word]]
    [{NEWLINE}
        [{"\textbackslash n"},tier=word]
    ]
    [{$<$block$>$}
        [{INDENT}[{$\rightarrow$},tier=word]]
        [{$<$assignment$>$}
            [{$<$identifier$>$}[{"x"},tier=word]]
            [{"="}[{"="},tier=word]]
            [{$<$expression$>$}
                [{$<$number$>$}[{"1"},tier=word]]
            ]
        ]
        [{DEDENT}[{$\leftarrow$},tier=word]]
    ]
]
\end{forest}
\end{center}

\vspace{5mm}

The derivation tree will produce the result:
\begin{lstlisting}
    for 2 times
        x = 1
\end{lstlisting}
The tree illustrates that the block scopes are handled through forced indentations and dedentations, a practice similar to \emph{Python}.
